\documentclass{article}
\usepackage{amsmath,amssymb,amsthm}
\usepackage{algorithm}
\usepackage{algorithmic}
\usepackage{enumitem}
\usepackage{hyperref}
\newtheorem{theorem}{Theorem}
\newtheorem{lemma}{Lemma}
\newcommand{\E}{\mathbb{E}}
\newcommand{\R}{\mathbb{R}}
\newcommand{\norm}[1]{\lVert #1\rVert}
\begin{document}

\section*{Randomized Block Coordinate Gradient Descent (RBCGD)}

\section*{Mathematical Formulation}
Let $f:\R^{d}\rightarrow\R$ be differentiable and written in a \emph{block} partition
\[
w = (w^{(1)},w^{(2)},\dots ,w^{(B)}),\qquad
w^{(i)}\in\R^{d_i},\;\;\sum_{i=1}^{B}d_i=d .
\]
At iteration $t\ge 0$ we draw a block index $i_t\in\{1,\dots ,B\}$ according to a
probability distribution $p_i>0$ (uniform $p_i=1/B$ is typical).  
Define the (partial) gradient w.r.t. block $i$ as
\[
\nabla_i f(w) \;:=\; \frac{\partial f(w)}{\partial w^{(i)}}\in\R^{d_i}.
\]

Given a stepsize schedule $\{\eta_t\}_{t\ge 0}$, the update reads
\[
\boxed{
\begin{aligned}
w^{(i_t)}_{t+1} &= w^{(i_t)}_{t} - \eta_t \,\nabla_{i_t} f(w_t),\\
w^{(j)}_{t+1} &= w^{(j)}_{t},\qquad\forall j\neq i_t .
\end{aligned}}
\tag{1}
\]

All other blocks remain unchanged.  The algorithm stops after a prescribed
maximum number of iterations $T$ or when a prescribed tolerance on the
gradient norm is reached.

\section*{Pseudocode}
\begin{algorithm}[H]
\caption{Randomized Block Coordinate Gradient Descent}
\begin{algorithmic}[1]
\STATE \textbf{Input:} initial point $w_0\in\R^{d}$, stepsize schedule $\{\eta_t\}$, 
maximum iterations $T$, block distribution $\{p_i\}_{i=1}^B$.
\FOR{$t=0,\dots ,T-1$}
    \STATE Sample block $i_t\sim\{p_i\}_{i=1}^B$.
    \STATE Compute partial gradient $g_t:=\nabla_{i_t}f(w_t)$.
    \STATE Update the selected block:
    \[
    w^{(i_t)}_{t+1}= w^{(i_t)}_{t}-\eta_t g_t .
    \]
    \STATE Set $w^{(j)}_{t+1}=w^{(j)}_{t}$ for all $j\neq i_t$.
\ENDFOR
\STATE \textbf{Output:} $\{w_t\}_{t=0}^{T}$ (or the last iterate $w_T$).
\end{algorithmic}
\end{algorithm}

\section*{Dry Run Example}
Consider the scalar function (single block) $f(w)=(w-3)^2$, $w\in\R$,
initial point $w_0=0$, constant stepsize $\eta_t=\eta=0.1$, and $T=3$.
Since there is only one block, $i_t=1$ for all $t$.

\begin{tabular}{c|c|c|c}
$t$ & $w_t$ & $\nabla f(w_t)=2(w_t-3)$ & $f(w_t)$\\ \hline
0 & $0$ & $-6$ & $9$\\
1 & $0-0.1(-6)=0.6$ & $2(0.6-3)=-4.8$ & $(0.6-3)^2=5.76$\\
2 & $0.6-0.1(-4.8)=1.08$ & $2(1.08-3)=-3.84$ & $(1.08-3)^2=3.6864$\\
3 & $1.08-0.1(-3.84)=1.464$ & $2(1.464-3)=-3.072$ & $(1.464-3)^2=2.3767$
\end{tabular}

The objective value decreases at each iteration, illustrating the
behaviour of (1).

\newpage

\section*{Assumptions}
\begin{enumerate}[label=(A\arabic*)]
\item \textbf{Lower Boundedness.}
\[
f^\star := \inf_{w\in\R^{d}} f(w) > -\infty .
\tag{A1}
\]
\item \textbf{Block-wise $L_i$-smoothness.}
For each block $i$ there exists $L_i>0$ such that for any $w,\tilde w$ that differ
only in block $i$,
\[
f(\tilde w) \le f(w) + \langle \nabla_i f(w),\tilde w^{(i)}-w^{(i)}\rangle
      + \frac{L_i}{2}\,\norm{\tilde w^{(i)}-w^{(i)}}^{2}.
\tag{A2}
\]
Define $L_{\max}:=\max_i L_i$.
\item \textbf{Unbiased Block Gradient Estimate.}
The random block selection yields an unbiased estimate of the full gradient:
\[
\E_{i_t}\big[\,\mathbf{e}_{i_t}\nabla_{i_t}f(w)\,\big]=\frac{1}{B}\nabla f(w),
\tag{A3}
\]
where $\mathbf{e}_{i}$ is the embedding operator that places a $d_i$‑vector
into the $d$‑dimensional space (zeros elsewhere).  Consequently
\[
\E_{i_t}\big[\norm{\mathbf{e}_{i_t}\nabla_{i_t}f(w)}^{2}\big]
      =\frac{1}{B}\,\norm{\nabla f(w)}^{2}.
\tag{A4}
\]
\item \textbf{Stepsize.}
The stepsize is constant $\eta_t\equiv\eta$ with
\[
0<\eta\le\frac{1}{L_{\max}} .
\tag{A5}
\]
\end{enumerate}

\section*{Main Convergence Theorem}
\begin{theorem}[Expected Gradient Norm Decay]\label{thm:main}
Under Assumptions (A1)–(A5) and for any $T\ge1$, the iterates generated by
\eqref{1} satisfy
\[
\frac{1}{T}\sum_{t=0}^{T-1}\E\!\left[\norm{\nabla f(w_t)}^{2}\right]
\;\le\;
\frac{2B\big(f(w_0)-f^\star\big)}{\eta\,T}.
\tag{2}
\]
Consequently,
\[
\min_{0\le t\le T-1}\E\!\left[\norm{\nabla f(w_t)}^{2}\right]
\;\le\;
\frac{2B\big(f(w_0)-f^\star\big)}{\eta\,T}
= O\!\left(\frac{1}{T}\right).
\]
\end{theorem}

\section*{Theoretical Properties}
\begin{enumerate}[label=\arabic*.]
\item \textbf{Stability.} The bound \eqref{2} requires only that
$\eta\le 1/L_{\max}$, guaranteeing that each block update does not increase
the objective in expectation.
\item \textbf{Hyper‑parameter Sensitivity.}
The rate constant is proportional to $B$; using larger blocks (smaller $B$)
improves the bound, reflecting the classic trade‑off between per‑iteration
cost and convergence speed.
\item \textbf{Comparison to Full‑gradient GD.}
When $B=1$ the algorithm reduces to standard (deterministic) gradient descent
and \eqref{2} recovers the well‑known $O(1/T)$ bound for non‑convex smooth
functions.  For $B>1$ the bound degrades by a factor $B$, matching the
information‑theoretic limit of block sampling.
\item \textbf{Special Cases.}
\begin{itemize}
\item \emph{Uniform sampling} $p_i=1/B$ yields the clean constants above.
\item \emph{Non‑uniform sampling} $p_i\propto L_i$ leads to a bound with
$L_{\rm avg}:=\frac{1}{B}\sum_i L_i$ replacing $L_{\max}$ in the stepsize
condition and a factor $\sum_i L_i/p_i$ in the numerator.
\end{itemize}
\end{enumerate}

\section*{Discussion}
The theorem shows that even for a non‑convex objective, the expected squared
norm of the full gradient diminishes at an $O(1/T)$ rate when only a random
block is updated each iteration.  This matches the best known rates for
stochastic gradient methods under comparable smoothness assumptions,
while enjoying a reduced per‑iteration computational burden.

\newpage

\section*{Preliminaries (Definitions and Lemmas)}
\begin{lemma}[Block Smoothness Inequality]\label{lem:smooth}
Let $i\in\{1,\dots ,B\}$ and $w,\tilde w$ differ only in block $i$.  Under (A2),
\[
f(\tilde w) \le f(w) - \eta\langle \nabla_i f(w),\nabla_i f(w)\rangle
          + \frac{L_i\eta^{2}}{2}\norm{\nabla_i f(w)}^{2}.
\tag{3}
\]
\end{lemma}
\begin{proof}
Apply (A2) with $\tilde w^{(i)} = w^{(i)} - \eta \nabla_i f(w)$ and
$\tilde w^{(j)}=w^{(j)}$ for $j\neq i$:
\[
\begin{aligned}
f(\tilde w)
&\le f(w) + \big\langle \nabla_i f(w), -\eta \nabla_i f(w)\big\rangle
     +\frac{L_i}{2}\norm{-\eta\nabla_i f(w)}^{2}\\
&= f(w) - \eta \norm{\nabla_i f(w)}^{2}
   +\frac{L_i\eta^{2}}{2}\norm{\nabla_i f(w)}^{2}.
\end{aligned}
\]
\end{proof}

\section*{Main Proof (Step‑by‑Step Derivation)}
We prove Theorem~\ref{thm:main}.  All expectations are taken w.r.t. the
random block $i_t$ conditioned on the past $\mathcal{F}_t:=\sigma(w_0,\dots ,w_t)$.

\noindent\textbf{[Step 1]} Apply Lemma~\ref{lem:smooth} with block $i_t$ and
$\eta_t\equiv\eta$:
\[
f(w_{t+1}) \le f(w_t) -\eta \norm{\nabla_{i_t}f(w_t)}^{2}
                 +\frac{L_{i_t}\eta^{2}}{2}\norm{\nabla_{i_t}f(w_t)}^{2}.
\tag{4}
\]

\noindent\textbf{[Step 2]} Use $L_{i_t}\le L_{\max}$ and factor:
\[
f(w_{t+1}) \le f(w_t) -\underbrace{\Bigl(\eta-\frac{L_{\max}\eta^{2}}{2}\Bigr)}_{=:c}
               \norm{\nabla_{i_t}f(w_t)}^{2}.
\tag{5}
\]

\noindent\textbf{[Step 3]} Take conditional expectation $\E[\cdot|\mathcal{F}_t]$.
By (A4),
\[
\E\big[\norm{\nabla_{i_t}f(w_t)}^{2}\mid\mathcal{F}_t\big]
   =\frac{1}{B}\norm{\nabla f(w_t)}^{2}.
\tag{6}
\]
Hence,
\[
\E\big[f(w_{t+1})\mid\mathcal{F}_t\big]
   \le f(w_t) - c\;\frac{1}{B}\norm{\nabla f(w_t)}^{2}.
\tag{7}
\]

\noindent\textbf{[Step 4]} Unroll the recursion by total expectation.
Define $\Delta_t:=\E\big[f(w_t)-f^\star\big]\ge0$.  From (7):
\[
\Delta_{t+1}\le \Delta_t - \frac{c}{B}\E\!\left[\norm{\nabla f(w_t)}^{2}\right].
\tag{8}
\]

\noindent\textbf{[Step 5]} Sum (8) from $t=0$ to $T-1$:
\[
\Delta_T \le \Delta_0 - \frac{c}{B}\sum_{t=0}^{T-1}
                \E\!\left[\norm{\nabla f(w_t)}^{2}\right].
\tag{9}
\]
Since $\Delta_T\ge0$, rearrange:
\[
\sum_{t=0}^{T-1}\E\!\left[\norm{\nabla f(w_t)}^{2}\right]
   \le \frac{B}{c}\,\Delta_0.
\tag{10}
\]

\noindent\textbf{[Step 6]} Choose $\eta$ as in (A5).  The smallest admissible
$\eta$ that maximises $c=\eta-\frac{L_{\max}\eta^{2}}{2}$ is $\eta=1/L_{\max}$,
giving $c=\frac{1}{2L_{\max}}$.  Substituting $c$ and $\Delta_0=f(w_0)-f^\star$:
\[
\sum_{t=0}^{T-1}\E\!\left[\norm{\nabla f(w_t)}^{2}\right]
   \le 2B L_{\max}\bigl(f(w_0)-f^\star\bigr).
\tag{11}
\]

\noindent\textbf{[Step 7]} Divide both sides of (11) by $T$ to obtain the
average bound:
\[
\frac{1}{T}\sum_{t=0}^{T-1}\E\!\left[\norm{\nabla f(w_t)}^{2}\right]
   \le \frac{2B L_{\max}\bigl(f(w_0)-f^\star\bigr)}{T}.
\tag{12}
\]
If we retain the original stepsize $\eta\le1/L_{\max}$, then $c\ge
\eta/2$, and (12) becomes the cleaner statement used in the theorem:
\[
\frac{1}{T}\sum_{t=0}^{T-1}\E\!\left[\norm{\nabla f(w_t)}^{2}\right]
   \le \frac{2B\bigl(f(w_0)-f^\star\bigr)}{\eta\,T}.
\tag{13}
\]
This completes the proof.  $\square$

\section*{Rate Derivation (With Constants)}
The bound \eqref{13} can be rewritten as
\[
\mathbb{E}\!\big[\|\nabla f(w_{\tau})\|^{2}\big]
   \le O\!\Bigl(\frac{B}{\eta T}\Bigr),
\qquad
\tau\sim\text{Uniform}\{0,\dots ,T-1\}.
\]
Thus, to achieve $\E[\|\nabla f(w_{\tau})\|^{2}]\le\epsilon$ we need
\[
T = O\!\Bigl(\frac{B}{\eta\epsilon}\Bigr).
\]

\section*{Special Cases}
\begin{itemize}
\item \textbf{Single block ($B=1$).} The algorithm coincides with
standard gradient descent and the bound reduces to
$\frac{2(f(w_0)-f^\star)}{\eta T}$.
\item \textbf{Full gradient stochastic method.} If each iteration samples
all blocks ($B=d$) but uses an unbiased stochastic gradient with variance
$\sigma^{2}$, an extra term $\frac{\eta\sigma^{2}}{B}$ appears in (12), leading
to the classical $O(1/\sqrt{T})$ rate for stochastic non‑convex optimization.
\item \textbf{Non‑uniform block probabilities.}
Let $p_i>0$, $\sum_i p_i=1$, and define $L_{\rm max}^{p}:=\max_i L_i/p_i$.
Choosing $\eta\le 1/L_{\rm max}^{p}$ yields the same proof with $B$ replaced
by $1/\min_i p_i$ in the numerator of \eqref{13}.
\end{itemize}

\end{document}